\section{Diseño y desarrollo de la solución de IA}\label{e-0}

Un asistente FinOps solo es útil si sus respuestas se pueden comprobar: una cifra sin origen o una recomendación que no se apoya en los datos del cliente no sirve para decidir gasto. Este apartado explica qué conocimiento prepara Economicon, cómo lo recupera, qué controles aplica antes de mostrar una respuesta y qué partes del diseño están activas en la configuración por defecto. La tabla~\ref{tab-e1} resume las decisiones; la evaluación de su calidad está en el apartado~\ref{h-0}.

\needspace{0.36\textheight}
\begin{table}[!h]
\centering\small
\caption{Decisiones de preparación, recuperación y generación}\label{tab-e1}
\begin{tabular}{@{}p{0.2\linewidth}p{0.3\linewidth}p{0.42\linewidth}@{}}
\toprule
\textbf{Elemento} & \textbf{Elección} & \textbf{Qué significa} \\
\midrule
Representación del texto & text-embedding-3-small; 1.536 dimensiones & Sirve para comparar cada pasaje con la pregunta. El código admite además Titan de AWS, apagado por defecto. \cite{especificaciones,config-gateway,arquitectura-final} \\
Fragmentación & Objetivo de 600 tokens; máximo de 800; solapamiento del 10\,\% & El corte respeta secciones, listas y tablas cuando caben. \cite{especificaciones} \\
Recuperación por defecto & Distancia coseno $\leq$ 0,6; hasta 4 pasajes & La distancia coseno mide cuánto se parecen dos vectores (0 es idéntico); solo se admiten pasajes por debajo del umbral. \cite{arquitectura-final} \\
Recuperación híbrida & Vectorial y por palabras, fusionadas por posición & Implementada y elegida en las campañas, pero apagada por defecto: requiere activarla y preparar un índice verificado. \cite{arquitectura-final,informe-rag} \\
Reintentos de una respuesta rechazada & Hasta 4 intentos; temperatura 0; 0,2; 0,4; 0,6 & La temperatura sube en cada reintento. Incluido en la versión final. \cite{cortes-medicion} \\
Consultas de costes & SQL determinista, sin modelo & Propiedad, gasto y anomalías se calculan en el servidor y se devuelven con su evidencia. \cite{arquitectura-final} \\
\bottomrule
\end{tabular}
\end{table}

\subsection{Un asistente con información de referencia}\label{e-1}

FinOps relaciona costes, decisiones y responsables. Un modelo de lenguaje puede explicar conceptos generales, pero no conoce los datos de un cliente ni las reglas de su organización. Economicon usa RAG: antes de redactar, selecciona información autorizada del cliente y la incorpora a la petición, de modo que la respuesta pueda citar sus fuentes en lugar de depender de la memoria del modelo. \cite{especificaciones}

El sistema separa dos tipos de pregunta. Las preguntas sobre costes ---cuánto gasta un proyecto, qué servicio ha crecido o por qué se marcó una anomalía--- se resuelven con consultas SQL en el servidor y devuelven las cifras exactas con su procedencia; ahí no interviene ningún modelo. Las preguntas abiertas pasan por el chat generativo, que recibe el resumen de costes verificado por el servidor y los pasajes documentales recuperados. La figura~\ref{fig-e1} muestra cómo se prepara el conocimiento y el ciclo de una respuesta. \cite{arquitectura-final}

El alcance cubre Azure con datos públicos de Microsoft FinOps Toolkit, servidos por una API que simula el servicio de consulta de costes; no se usa la facturación de un cliente real. Cada organización, o tenant, solo accede a sus documentos, costes y conversaciones. \cite{especificaciones,dataset-azure}

\figura{ia-flujo-conocimiento-y-respuesta}{Preparación del conocimiento y ciclo de una respuesta.\label{fig-e1}}

\subsection{Preparación del conocimiento}\label{e-2}

La ingesta admite texto pegado y ficheros Markdown, TXT, DOCX, DOC y PDF, incluidos los escaneados. La extracción y el OCR se hacen dentro del sistema, aislados y con límites de páginas, tiempo y memoria. El fichero original se guarda de forma temporal en MinIO y se retira cuando el texto normalizado, los vectores y el estado del trabajo ya están guardados. Así, la recuperación y las citas no dependen de conservar el binario. \cite{arquitectura-final}

Un fragmento (chunk) es la unidad de texto que se guarda y se recupera por separado. El algoritmo conserva secciones, listas, tablas y bloques de código cuando caben en el presupuesto de la tabla~\ref{tab-e1} y aplica cortes reproducibles cuando no caben, sin usar modelos para decidir dónde cortar. El solapamiento mantiene la continuidad entre fragmentos y queda registrado, de forma que el texto normalizado se puede reconstruir. \cite{especificaciones}

\subsection{Representación y recuperación}\label{e-3}

PostgreSQL con la extensión pgvector actúa como base vectorial: guarda los embeddings y devuelve los pasajes más próximos a la pregunta, siempre filtrados por tenant. Documentos y preguntas deben compartir el mismo espacio de representación ---modelo, proveedor, dimensión y preparación del texto---; cambiar de proveedor exige reindexar, nunca mezclar vectores. \cite{arquitectura-final}

Para elegir la recuperación se compararon cuatro variantes acumulativas: A usaba fragmentos por caracteres y búsqueda solo vectorial; B introdujo la fragmentación estructural; C añadió título y sección a la representación; D sumó la búsqueda por palabras, útil cuando importan un nombre, una magnitud o un identificador literal. D combina ambas listas con Reciprocal Rank Fusion (RRF), que puntúa la posición de cada candidato sin mezclar escalas distintas. \cite{informe-rag}

Elegimos D porque obtuvo respuesta en todas las preguntas de su campaña y seleccionó contexto más útil en las preguntas específicas; A y B cubrieron mejor algunas de las originales. Es una decisión ligada a ese corpus y a esa muestra, no una superioridad general. En la versión final, la recuperación verificada, contextual e híbrida está implementada pero apagada por defecto: activarla exige preparar el índice por generaciones completas, y mientras tanto el sistema usa la recuperación vectorial con umbral de la tabla~\ref{tab-e1}. \cite{informe-rag,arquitectura-final}

\subsection{Generación y controles}\label{e-4}

El chat generativo también se activa de forma explícita; en la configuración por defecto la interfaz no lo muestra. Cuando está activo, el modelo no decide si se busca ni qué se busca: el servidor prepara una única búsqueda por intento a partir de la pregunta y la ejecuta antes del primer borrador. El modelo recibe solo ese resultado y el resumen de costes del cliente; no ejecuta consultas ni accede a otros datos. \cite{arquitectura-final}

La respuesta sigue un esquema estructurado ---campos con tipos y límites definidos--- y un prompt versionado. Antes de guardarla, el servidor comprueba que cada cifra, periodo y cita se apoya en la evidencia autorizada. Si una cifra no se puede verificar, se omite con un aviso visible y se conserva el resto de la respuesta; si la salida no es válida, se reintenta con la variabilidad creciente de la tabla~\ref{tab-e1}. La conversación se recupera tras un fallo o una cancelación, y el historial largo se resume para acotar el contexto que recibe el modelo. \cite{cortes-medicion,especificaciones}

Las consultas de propiedad, gasto y anomalías todavía no están disponibles como herramientas que el modelo pueda invocar: se piden mediante formularios explícitos. Conectarlas al chat es el siguiente paso del diseño (apartado~\ref{i-0}). \cite{arquitectura-final}
